\section{Indsætte kodeeksempler}
\subsection{Java kode}
Det er muligt at specificere sproget, som man skriver i ved brug af "Minted":
\begin{minted}{java}
/**
 * Vi har har startet et minted environment:
 * - Ved brug af minted package kan man fremhæve tekst
 * - Så  det ligner den klassiske IDE
 * - Vi har et eksempel her med assert statements
 * - Der bliver tilføjet en cake, hvis den ikke allerede eksisterer
 */
public void addCake(Cake c) throws DuplicateException {
    assert cakes != null;
    assert c != null;
		
    if(containsCakeWithName(c.getName())) {
        assert containsCakeWithName(c.getName());
        throw new DuplicateException("Duplicate cake in kitchen " + this.getName());
    }
    this.cakes.add(c);
    assert containsCakeWithName(c.getName());
}
\end{minted}

\subsection{Python kode}
Det er også muligt at ændre baggrundsfarven!
\usemintedstyle{monokai}
\begin{minted}
[
frame=lines,
framesep=2mm,
baselinestretch=1.2,
bgcolor=black,
fontsize=\footnotesize,
linenos
]
{python}
import numpy as np
    
def incmatrix(genl1,genl2):
    m = len(genl1)
    n = len(genl2)
    M = None #to become the incidence matrix
    VT = np.zeros((n*m,1), int)  #dummy variable
    
    #compute the bitwise xor matrix
    M1 = bitxormatrix(genl1)
    M2 = np.triu(bitxormatrix(genl2),1) 

    for i in range(m-1):
        for j in range(i+1, m):
            [r,c] = np.where(M2 == M1[i,j])
            for k in range(len(r)):
                VT[(i)*n + r[k]] = 1;
                VT[(i)*n + c[k]] = 1;
                VT[(j)*n + r[k]] = 1;
                VT[(j)*n + c[k]] = 1;
                
                if M is None:
                    M = np.copy(VT)
                else:
                    M = np.concatenate((M, VT), 1)
                
                VT = np.zeros((n*m,1), int)
    
    return M
\end{minted}

\subsection{MatLab kode}
Man kan også blot bruge et lstlisting environment og ikke specificere sproget:
\begin{lstlisting}
% For at indsaette kode fra Matlab, kan man blot kopiere sit script og indsaette koden her
\end{lstlisting}
Eksempel på en meget fjollet kode, som tjekker den alder, du indtaster. Hvis man kører koden, spørger den først "Please enter your name", herefter spørger den om din alder. Hvis du enten taster et tal mindre end 18 eller bogstaver, får du en fejlmeddelelse. 
\begin{lstlisting}
%% Please tell your mom
name = input('Please enter your name: ', 's');
while true
age = str2double(input('Enter your age in years: ', 's'));
if ~isnan(age)  && (age > 18) % needs to be over 18
    
break;
end
if age < 18 
    disp('You are not old enough. Please tell your mom.')
else 
    disp('Not a valid number. Please try again.')
end
end

\end{lstlisting}
